
\documentclass{article}
\usepackage{colortbl}
\usepackage{makecell}
\usepackage{multirow}
\usepackage{supertabular}

\begin{document}

\newcounter{utterance}

\centering \large Interaction Transcript for game `hot\_air\_balloon', experiment `air\_balloon\_survival\_mt\_complexity\_hard', episode 2 with gemma{-}4{-}26b{-}a4b{-}it.
\vspace{24pt}

{ \footnotesize  \setcounter{utterance}{1}
\setlength{\tabcolsep}{0pt}
\begin{supertabular}{c@{$\;$}|p{.15\linewidth}@{}p{.15\linewidth}p{.15\linewidth}p{.15\linewidth}p{.15\linewidth}p{.15\linewidth}}
    \# & $\;$A & \multicolumn{4}{c}{Game Master} & $\;\:$B\\
    \hline

    \theutterance \stepcounter{utterance}  
    & & \multicolumn{4}{p{0.6\linewidth}}{
        \cellcolor[rgb]{0.9,0.9,0.9}{
            \makecell[{{p{\linewidth}}}]{
                \texttt{\tiny{[P1$\langle$GM]}}
                \texttt{Int qiegħed tipparteċipa f'logħba ta' negozjar kollaborattiva.} \\
\\ 
\texttt{Flimkien ma' parteċipant ieħor, trid taqblu fuq sett wieħed ta' oġġetti li se jinżammu. Kull wieħed minnkom għandu l{-}fehma tiegħu dwar kemm kull oġġett hu importanti għalih (importanza). Int ma tafx kif il{-}parteċipant l{-}ieħor jivvaluta l{-}oġġetti. Barra minn hekk, tingħata l{-}isforz li jitlob kull oġġett.} \\
\texttt{Tistgħu taqblu fuq sett biss jekk l{-}isforz totali tal{-}oġġetti magħżula ma jaqbiżx limitu kondiviż:} \\
\\ 
\texttt{LIMITU = 4843} \\
\\ 
\texttt{Hawn huma l{-}valuri individwali tal{-}isforz tal{-}oġġetti:} \\
\\ 
\texttt{Isforz tal{-}oġġett = \{"C08": 880, "C55": 158, "A43": 104, "A67": 514, "B42": 797, "A62": 815, "B19": 336, "C43": 79, "C73": 522, "C10": 973, "B36": 682, "A18": 178, "C61": 184, "C67": 795, "C42": 154, "C70": 145, "B40": 842, "C39": 887, "A21": 328, "C16": 313\}} \\
\\ 
\texttt{Hawn hi l{-}fehma personali tiegħek dwar l{-}importanza ta' kull oġġett:} \\
\\ 
\texttt{Valuri tal{-}importanza tal{-}oġġetti = \{"C08": 977, "C55": 255, "A43": 201, "A67": 611, "B42": 894, "A62": 912, "B19": 433, "C43": 176, "C73": 619, "C10": 1070, "B36": 779, "A18": 275, "C61": 281, "C67": 892, "C42": 251, "C70": 242, "B40": 939, "C39": 984, "A21": 425, "C16": 410\}} \\
\\ 
\texttt{Għan:} \\
\\ 
\texttt{L{-}għan tiegħek hu li tinnegozja sett kondiviż ta' oġġetti li jibbenefikak kemm jista' jkun (jiġifieri, jimmassimizza l{-}importanza totali GĦALIK), filwaqt li tibqa' fil{-}LIMITU. M'intix obbligat tagħmel PROPOSTA f'kull messaġġ {-} tista' sempliċement tinnegozja wkoll. It{-}tattiċi kollha huma permessi!} \\
\\ 
\texttt{Protokoll ta' Interazzjoni:} \\
\\ 
\texttt{Tista' tuża biss il{-}formati strutturati li ġejjin f'messaġġ:} \\
\\ 
\texttt{PROPOSTA: \{'A', 'B', 'C', …\}} \\
\texttt{Ipproponi li jinżammu eżattament dawk l{-}oġġetti.} \\
\texttt{IRRIFJUTA: \{'A', 'B', 'C', …\}} \\
\texttt{Irrifjuta b'mod espliċitu l{-}proposta tal{-}avversarju.} \\
\texttt{ARGUMENT: \{'...'\}} \\
\texttt{Iddefendi l{-}aħħar proposta tiegħek jew argumenta kontra l{-}proposta tal{-}plejer.} \\
\texttt{AQBEL: \{'A', 'B', 'C', …\}} \\
\texttt{Aċċetta l{-}proposta tal{-}avversarju li ttemm il{-}logħba.} \\
\\ 
\\ 
\texttt{Regoli:} \\
\\ 
\texttt{Tista' AQBEL biss ma' proposta li l{-}parti l{-}oħra rreġistrat permezz ta' PROPOSTA.} \\
\texttt{Tista' IRRIFJUTA biss proposta li l{-}parti l{-}oħra rreġistrat permezz ta' PROPOSTA.} \\
\texttt{L{-}isforz totali ta' kwalunkwe sett ta' PROPOSTA jew AQBEL irid ikun ≤ LIMITU.} \\
\texttt{TIŻVELAX il{-}punteġġi moħbija tal{-}importanza tiegħek.} \\
\texttt{Tikketta f'format strutturat trid tkun segwita minn kolon u spazju vojt. L{-}argument irid ikun sett ta' python li fih 0 jew aktar strings.} \\
\texttt{Għalhekk, irid ikun fil{-}forma "TIKKETTA: \{...\}"} \\
\texttt{Segwi b'mod strett il{-}protokoll ta' interazzjoni u TIKTIBX xejn lil hinn mill{-}istruttura mogħtija.} \\
\texttt{Il{-}logħba tispiċċa meta naħa waħda tagħti AQBEL għal PROPOSTA magħmula mill{-}plejer l{-}ieħor.} \\
\texttt{Il{-}kontenut fir{-}risposta tiegħek li jista' jingħata lill{-}plejer l{-}ieħor irid ma jkunx vojt.} \\
\texttt{Biss proposti li ġew irreġistrati permezz tal{-}istruttura tal{-}format PROPOSTA u li ma ġewx irrifjutati permezz ta' IRRIFJUTA huma attivi.} \\
\\ 
            }
        }
    }
    & \\ \\

    \theutterance \stepcounter{utterance}  
    & \multicolumn{4}{p{0.6\linewidth}}{
        \cellcolor[rgb]{0.9,0.9,0.9}{
            \makecell[{{p{\linewidth}}}]{
                \texttt{\tiny{[P1$\rangle$GM]}}
                \texttt{PROPOSTA: \{'C08', 'C55', 'A43', 'A67', 'B42', 'A62', 'B19', 'C43', 'C73', 'C10', 'B36', 'A18', 'C61', 'C67', 'C42', 'C70', 'B40', 'C39', 'A21', 'C16'\}} \\
\\ 
\texttt{ARGUMENT: \{'Bdejna b'proposta li tinkludi l{-}oġġetti kollha biex nittestjaw l{-}isforz totali u l{-}importanza. L{-}isforz totali għal dawn l{-}oġġetti kollha huwa 10544, li huwa ħafna fuq il{-}limit tard. Għandna nibdew nipproponu settijiet li jiffokaw fuq l{-}oġġetti li għandhom valur għal qedna b'mod konċiż.'\}} \\
            }
        }
    }
    & & \\ \\

    \theutterance \stepcounter{utterance}  
    & & \multicolumn{4}{p{0.6\linewidth}}{
        \cellcolor[rgb]{0.9,0.9,0.9}{
            \makecell[{{p{\linewidth}}}]{
                \texttt{\tiny{[GM$\rangle$P2]}}
                \texttt{Qed tipparteċipa f'logħba ta' negozjar kollaborattiva.} \\
\\ 
\texttt{Flimkien ma' parteċipant ieħor, trid taqblu fuq sett wieħed ta' oġġetti li se jinżammu. Kull wieħed minnkom għandu l{-}fehma tiegħu dwar kemm kull oġġett hu importanti għalih (importanza). Ma tafx kif il{-}parteċipant l{-}ieħor jivvaluta l{-}oġġetti. Barra minn hekk, tingħata l{-}isforz li jitlob kull oġġett.} \\
\texttt{Tistgħu taqblu fuq sett biss jekk l{-}isforz totali tal{-}oġġetti magħżula ma jaqbiżx limitu kondiviż:} \\
\\ 
\texttt{LIMITU = 4843} \\
\\ 
\texttt{Hawn huma l{-}valuri individwali tal{-}isforz tal{-}oġġetti:} \\
\\ 
\texttt{Isforz tal{-}oġġett = \{"C08": 880, "C55": 158, "A43": 104, "A67": 514, "B42": 797, "A62": 815, "B19": 336, "C43": 79, "C73": 522, "C10": 973, "B36": 682, "A18": 178, "C61": 184, "C67": 795, "C42": 154, "C70": 145, "B40": 842, "C39": 887, "A21": 328, "C16": 313\}} \\
\\ 
\texttt{Hawn hi l{-}fehma personali tiegħek dwar l{-}importanza ta' kull oġġett:} \\
\\ 
\texttt{Valuri tal{-}importanza tal{-}oġġetti = \{"C08": 977, "C55": 255, "A43": 201, "A67": 611, "B42": 894, "A62": 912, "B19": 433, "C43": 176, "C73": 619, "C10": 1070, "B36": 779, "A18": 275, "C61": 281, "C67": 892, "C42": 251, "C70": 242, "B40": 939, "C39": 984, "A21": 425, "C16": 410\}} \\
\\ 
\texttt{Għan:} \\
\\ 
\texttt{L{-}għan tiegħek hu li tinnegozja sett kondiviż ta' oġġetti li jibbenefikak kemm jista' jkun (jiġifieri, jimmassimizza l{-}importanza totali GĦALIK), filwaqt li tibqa' fil{-}LIMITU. M'intix obbligat tagħmel PROPOSTA f'kull messaġġ {-} tista' sempliċement tinnegozja wkoll. It{-}tattiċi kollha huma permessi!} \\
\\ 
\texttt{Protokoll ta' Interazzjoni:} \\
\\ 
\texttt{Tista' tuża biss il{-}formati strutturati li ġejjin f'messaġġ:} \\
\\ 
\texttt{PROPOSTA: \{'A', 'B', 'C', …\}} \\
\texttt{Ipproponi li jinżammu eżatt dawk l{-}oġġetti.} \\
\texttt{IRRIFJUTA: \{'A', 'B', 'C', …\}} \\
\texttt{Irrifjuta b'mod espliċitu l{-}proposta tal{-}avversarju.} \\
\texttt{ARGUMENT: \{'...'\}} \\
\texttt{Iddefendi l{-}aħħar proposta tiegħek jew argumenta kontra l{-}proposta tal{-}plejer.} \\
\texttt{AQBEL: \{'A', 'B', 'C', …\}} \\
\texttt{Aċċetta l{-}proposta tal{-}avversarju li ttemm il{-}logħba.} \\
\\ 
\\ 
\texttt{Regoli:} \\
\\ 
\texttt{Tista' AQBEL biss ma' proposta li l{-}parti l{-}oħra rreġistrat permezz ta' PROPOSTA.} \\
\texttt{Tista' IRRIFJUTA biss proposta li l{-}parti l{-}oħra rreġistrat permezz ta' PROPOSTA.} \\
\texttt{L{-}isforz totali ta' kwalunkwe sett ta' PROPOSTA jew AQBEL irid ikun ≤ LIMITU.} \\
\texttt{TIŻVELAX il{-}punteġġi moħbija tal{-}importanza tiegħek.} \\
\texttt{Tikketta f'format strutturat trid tkun segwita b'kolon u spazju. L{-}argument irid ikun sett ta' Python li jkun fih 0 jew aktar strings.} \\
\texttt{Għalhekk, irid ikun fil{-}forma "TIKKETTA: \{...\}"} \\
\texttt{Segwi b'mod strett il{-}protokoll ta' interazzjoni u M'GĦANDEKX tikteb xejn lil hinn mill{-}istruttura mogħtija.} \\
\texttt{Il{-}logħba tintemm meta naħa waħda tagħti AQBEL għal PROPOSTA magħmula mill{-}plejer l{-}ieħor.} \\
\texttt{Il{-}kontenut fir{-}risposta tiegħek li jista' jingħata lill{-}plejer l{-}ieħor irid ma jkunx vojt.} \\
\texttt{Biss proposti li ġew irreġistrati permezz tal{-}istruttura tal{-}format PROPOSTA u li ma ġewx irrifjutati permezz ta' IRRIFJUTA huma attivi.} \\
\\ 
\\ 
\\ 
\texttt{Issa se tirċievi l{-}ewwel messaġġ tal{-}plejer l{-}ieħor.} \\
\\ 
\texttt{PROPOSTA: \{'C08', 'C55', 'A43', 'A67', 'B42', 'A62', 'B19', 'C43', 'C73', 'C10', 'B36', 'A18', 'C61', 'C67', 'C42', 'C70', 'B40', 'C39', 'A21', 'C16'\}} \\
\\ 
\texttt{ARGUMENT: \{'Bdejna b'proposta li tinkludi l{-}oġġetti kollha biex nittestjaw l{-}isforz totali u l{-}importanza. L{-}isforz totali għal dawn l{-}oġġetti kollha huwa 10544, li huwa ħafna fuq il{-}limit tard. Għandna nibdew nipproponu settijiet li jiffokaw fuq l{-}oġġetti li għandhom valur għal qedna b'mod konċiż.'\}} \\
            }
        }
    }
    & \\ \\

    \theutterance \stepcounter{utterance}  
    & & & \multicolumn{4}{p{0.6\linewidth}}{
        \cellcolor[rgb]{0.9,0.9,0.9}{
            \makecell[{{p{\linewidth}}}]{
                \texttt{\tiny{[GM$\langle$P2]}}
                \texttt{PROPOSTA: \{'C08', 'C10', 'B42', 'A62', 'C39', 'B40', 'C67'\}} \\
\\ 
\texttt{ARGUMENT: \{'L{-}isforz totali ta' din il{-}proposta huwa 5388, li għadu jaqbiż fuq il{-}limit. Għandna napprossimaw sett li jiffoka fuq l{-}oġġetti li għandhom l{-}ogħla importanza għalina b'mod effiċjenti. Din il{-}proposta hija l{-}ewwel pass biex nidentifikaw l{-}oġġetti li għandhom l{-}aktar valur.'\}} \\
            }
        }
    }
     \\ \\

    \theutterance \stepcounter{utterance}  
    & & \multicolumn{4}{p{0.6\linewidth}}{
        \cellcolor[rgb]{0.9,0.9,0.9}{
            \makecell[{{p{\linewidth}}}]{
                \texttt{\tiny{[P1$\langle$GM]}}
                \texttt{PROPOSTA: \{'C08', 'C10', 'B42', 'A62', 'C39', 'B40', 'C67'\}} \\
\\ 
\texttt{ARGUMENT: \{'L{-}isforz totali ta' din il{-}proposta huwa 5388, li għadu jaqbiż fuq il{-}limit. Għandna napprossimaw sett li jiffoka fuq l{-}oġġetti li għandhom l{-}ogħla importanza għalina b'mod effiċjenti. Din il{-}proposta hija l{-}ewwel pass biex nidentifikaw l{-}oġġetti li għandhom l{-}aktar valur.'\}} \\
            }
        }
    }
    & \\ \\

    \theutterance \stepcounter{utterance}  
    & \multicolumn{4}{p{0.6\linewidth}}{
        \cellcolor[rgb]{0.9,0.9,0.9}{
            \makecell[{{p{\linewidth}}}]{
                \texttt{\tiny{[P1$\rangle$GM]}}
                \texttt{IRRIFJUTA: \{'C08', 'C10', 'B42', 'A62', 'C39', 'B40', 'C67'\}} \\
\\ 
\texttt{ARGUMENT: \{'L{-}isforz totali ta' din il{-}proposta huwa 5388, li huwa għal kollox fuq il{-}limitu tagħna ta' 4843. Għandna nipproponu sett li jkun taħt il{-}limit.'\}} \\
            }
        }
    }
    & & \\ \\

    \theutterance \stepcounter{utterance}  
    & & \multicolumn{4}{p{0.6\linewidth}}{
        \cellcolor[rgb]{0.9,0.9,0.9}{
            \makecell[{{p{\linewidth}}}]{
                \texttt{\tiny{[GM$\rangle$P2]}}
                \texttt{IRRIFJUTA: \{'C08', 'C10', 'B42', 'A62', 'C39', 'B40', 'C67'\}} \\
\\ 
\texttt{ARGUMENT: \{'L{-}isforz totali ta' din il{-}proposta huwa 5388, li huwa għal kollox fuq il{-}limitu tagħna ta' 4843. Għandna nipproponu sett li jkun taħt il{-}limit.'\}} \\
            }
        }
    }
    & \\ \\

    \theutterance \stepcounter{utterance}  
    & & & \multicolumn{4}{p{0.6\linewidth}}{
        \cellcolor[rgb]{0.9,0.9,0.9}{
            \makecell[{{p{\linewidth}}}]{
                \texttt{\tiny{[GM$\langle$P2]}}
                \texttt{PROPOSTA: \{'C08', 'C10', 'B42', 'A62', 'C39', 'C67'\}} \\
\\ 
\texttt{ARGUMENT: \{'Għandek raġun. Din il{-}proposta l{-}ġdida tnaqqas l{-}isforz għal 5388 {-}> 4508, li huwa taħt il{-}limitu ta' 4843. Din hija proposta aktar realistika biex nibdew negozjar.'\}} \\
            }
        }
    }
     \\ \\

    \theutterance \stepcounter{utterance}  
    & & \multicolumn{4}{p{0.6\linewidth}}{
        \cellcolor[rgb]{0.9,0.9,0.9}{
            \makecell[{{p{\linewidth}}}]{
                \texttt{\tiny{[P1$\langle$GM]}}
                \texttt{PROPOSTA: \{'C08', 'C10', 'B42', 'A62', 'C39', 'C67'\}} \\
\\ 
\texttt{ARGUMENT: \{'Għandek raġun. Din il{-}proposta l{-}ġdida tnaqqas l{-}isforz għal 5388 {-}> 4508, li huwa taħt il{-}limitu ta' 4843. Din hija proposta aktar realistika biex nibdew negozjar.'\}} \\
            }
        }
    }
    & \\ \\

    \theutterance \stepcounter{utterance}  
    & \multicolumn{4}{p{0.6\linewidth}}{
        \cellcolor[rgb]{0.9,0.9,0.9}{
            \makecell[{{p{\linewidth}}}]{
                \texttt{\tiny{[P1$\rangle$GM]}}
                \texttt{AQBEL: \{'C08', 'C10', 'B42', 'A62', 'C39', 'C67'\}} \\
            }
        }
    }
    & & \\ \\

    \theutterance \stepcounter{utterance}  
    & & & \multicolumn{2}{p{0.3\linewidth}}{
        \cellcolor[rgb]{0.9,0.9,0.9}{
            \makecell[{{p{\linewidth}}}]{
                \texttt{\tiny{[GM$|$GM]}}
                \texttt{game successful} \\
            }
        }
    }
    & & \\ \\

    \theutterance \stepcounter{utterance}  
    & & & \multicolumn{2}{p{0.3\linewidth}}{
        \cellcolor[rgb]{0.9,0.9,0.9}{
            \makecell[{{p{\linewidth}}}]{
                \texttt{\tiny{[GM$|$GM]}}
                \texttt{end game} \\
            }
        }
    }
    & & \\ \\

\end{supertabular}
}

\end{document}
